% ======================================================================
%  KW06 中译本 · 思路导读（独立文档，与译文主文件互不干扰）
%  编译：在本目录下 xelatex 两遍；文件名 main.tex
% ======================================================================
\documentclass[11pt,a4paper]{article}
\usepackage[UTF8,fontset=windows,heading=true]{ctex}
\usepackage{amsmath,amssymb,amsthm}
\usepackage[margin=2.5cm]{geometry}
\usepackage{booktabs,array,longtable}
\usepackage{xcolor}
\usepackage[colorlinks,linkcolor=blue!50!black,urlcolor=blue!50!black]{hyperref}
\usepackage{enumitem}
\setlist{itemsep=2pt,topsep=2pt}
\newcommand{\MH}{\mathcal{M}_H}
\newcommand{\Ldual}{{}^L\!G}
\newcommand{\LdualC}{{}^L\!G_{\mathbb{C}}}
\newcommand{\Nfour}{\mathcal{N}=4}
\newcommand{\half}{\tfrac12}
\newcommand{\src}[1]{{\color{gray!70}\footnotesize[#1]}}
\newcommand{\kw}{\textsf{KW06}}
\linespread{1.15}

\title{\bfseries Kapustin--Witten (2006) 中译本\\[4pt]\large 思路导读}
\author{译本附册}
\date{2026 年 11 月}

\begin{document}
\maketitle
\thispagestyle{empty}

\begin{center}
\begin{minipage}{0.92\textwidth}
\small\itshape
这份导读回答一个问题：\kw{} 这 231 页，作者到底想干什么，读者应该按什么顺序读，
在哪里会卡住、卡住时该看什么。它不是摘要，也不是证明，而是\textbf{读之前的路线图}。
正文里的页锚一律给\textbf{印刷页}（= PDF 页 $-\,1$），与本译本边注一致。
\end{minipage}
\end{center}

\vspace{1em}

\section*{一、一句话与三步走}

\textbf{一句话}：Kapustin 与 Witten 把四维 $\Nfour$ 超对称规范理论的电磁对偶当成一台机器，
用它在黎曼面上生产出几何 Langlands 纲领所需要的 Hecke 特征层，并且诚实地说清楚了：
这台机器在物理上高度可信，但在数学上还不是证明。

由此可以给出三种读法：

\begin{enumerate}[label=(\arabic*)]
  \item \textbf{两小时版}（想知道结论）：读 \S1 引言、\S5.4--5.6（$S$-对偶 $=$ 镜像对称）、
        \S11.1（典范余各向同性 A-膜）。这三块是全篇的三个齿轮。
  \item \textbf{一天版}（想跟上推导）：再加 \S2（钉死对偶群）、\S3（扭化）、\S4（紧化与 $\MH$）、
        \S6（线算子）、\S9（'t Hooft $=$ Hecke）。跳过 \S7 的离散通量与 \S10 的技术细节。
  \item \textbf{两周版}（想吃透）：按 \S\,2 $\to$ 3 $\to$ 4 $\to$ 5 $\to$ 6 $\to$ 8 $\to$ 9 $\to$ 10 $\to$ 11 $\to$ 12
        逐节推进，每节读完回来对照本导读第三节的「卡点」。
\end{enumerate}

\section*{二、十二节地图}

下表是本译本与原文的对照。\textbf{任务}一栏写"这节要干什么"，\textbf{产物}一栏写"读完手里多了什么"，
\textbf{第一遍}一栏是建议：\textbf{必读}的节构成主链，\textbf{可略}的节第一遍可以先跳过。

\begin{center}\small
\begin{longtable}{@{}p{0.6cm}p{3.4cm}p{1.1cm}p{4.6cm}p{4.6cm}@{}}
\toprule
\# & 原文节（中译名） & 印刷页 & 任务 & 产物 / 卡点 \\
\midrule
\endhead
1 & 引言 & 2--5 & 交代六个基本想法 & 全篇目录；读不懂很正常，读完 \S5 再回来看 \\
2 & $\Nfour$ 超对称 Yang--Mills 与 $S$-对偶 & 6--15 & 钉死"对偶交换的是哪个群" & 词表：$G$、$\Ldual$、电荷格、复耦合 $\tau$ \\
3 & 由 $\Nfour$ SYM 得到拓扑场论 & 16--37 & 扭化，得到一族四维 TFT & 典范参数 $\Psi$——全文的旋钮 \\
4 & 紧化与 $\MH${} 的几何 & 38--55 & 紧化到曲线 $C$ 上，舞台出现 & 超凯勒流形 $\MH$，三复结构 $I,J,K$，Hitchin 纤维化 \\
5 & 二维拓扑场论 & 56--77 & $S$-对偶在二维变成镜像对称 & $B(J)\cong A(K)$；广义复几何。\textbf{全篇第一个核心} \\
6 & 圈算子与线算子 & 78--102 & 线算子变成膜上的算子 & 特征膜（eigenbrane）的定义 \\
7 & 通量与 $S$-对偶 & 103--115 & 群中心的离散通量 & 相位 $\pi_1(G)$ 的精细修正；可略 \\
8 & 电特征膜 & 116--122 & 零膜作电特征膜 & 特征膜 $\leftrightarrow$ $\pi_1(C)\to\Ldual$ 的表示 \\
9 & 't Hooft 算子与 Hecke 算子 & 123--147 & 两者是同一件事 & Hecke 修改的同构类、仿射 Grassmann。\textbf{第二个核心} \\
10 & Bogomolny 方程与 Hecke 修改 & 148--183 & 用单极子冒泡把 \S9 算出来 & 边界条件、Kähler 商、OPE、扩展 Bogomolny 方程 \\
11 & A-膜与 $\mathcal{D}$-模 & 184--206 & 把 A-膜翻译成 $\mathcal{D}$-模 & c.c.\ A-膜、开弦 BRST 代数。\textbf{收口} \\
12 & 由规范理论得到膜 & 207--222 & 把膜提升回四维规范理论 & 三类膜的边界条件；\textbf{诚实度指标} \\
\bottomrule
\end{longtable}
\end{center}

\section*{三、五条主线（五个"不显然之处"）}

全篇的骨头是五个 miracle。每个 miracle 都是一句"如果这成立，对应就成立"。

\textbf{M1｜扭化后仍是拓扑场论，且只依赖一个参数}。%
四维 $\Nfour$ 的超对称代数允许一族扭化，其中有一个复参数 $\Psi$（canonical parameter）。
关键事实是：整族理论都保持拓扑性，并且 $\Psi$ 完全刻画它——它不是耦合常数，而是理论的「身份证」。
$\Psi=\pm i$ 是退化点，此时理论回到通常的「几何 Langlands 参数」。\src{\S3，源 pp.16--37}

\smallskip
\textbf{M2｜四维的 $S$-对偶在二维恰好变成镜像对称}。%
把 $\Nfour$ 紧化在曲线 $C$ 上，得到 $\MH$ 上的二维 $\sigma$-模型。在 \S5 中作者论证：
四维的 $S$-对偶在这个二维理论里表现为复结构 $J\leftrightarrow K$ 的交换，也就是镜像对称。
这是全文最漂亮的一步——需要的对偶不是另找的，就是已有的那个。\src{\S5，源 pp.56--77}

\smallskip
\textbf{M3｜'t Hooft 算子的作用 $=$ Hecke 修改}。%
在 \S9 中作者把 't Hooft 线算子沿 $C$ 上的点插入，发现它对 Hitchin 模空间上的膜的作用，
与几何 Langlands 里 Hecke 算子的作用完全一致。\S10 再用单极子冒泡（monopole bubbling）
把这件事实算出来：单极子解的空间 $\leftrightarrow$ Hecke 修改的模空间。\src{\S9--\S10，源 pp.123--183}

\smallskip
\textbf{M4｜Wilson 线的特征膜给出局部系统}。%
电的 Wilson 线的特征膜，恰好对应 $\pi_1(C)\to\LdualC$ 的表示，即局部系统（平坦联络）。
这把几何 Langlands 的「自守侧」变成物理对象。\src{\S8，源 pp.116--122}

\smallskip
\textbf{M5｜典范余各向同性 A-膜把 A-膜范畴送到 $\mathcal{D}$-模范畴}。%
在 \S11 中，作者引入典范余各向同性 A-膜 $\mathcal{B}_{cc}$，并论证
$\mathcal{D}$-模可以被看成 $\mathcal{B}_{cc}$ 与别的膜之间的空间。翻译器因此闭合：
膜（物理）$\longleftrightarrow$ $\mathcal{D}$-模（几何）。\src{\S11，源 pp.184--206}

最短的理解路径可以压缩成一条链：
\[
\MH \;\longrightarrow\; \text{膜} \;\longrightarrow\; \text{线算子的本征性}
\;\longrightarrow\; \text{Hecke 算子} \;\longrightarrow\; \mathcal{D}\text{-模}.
\]
四个箭头上分别站着 M2、M4、M3、M5。其余所有概念都是这条链的解释或技术支撑。

\section*{四、读的时候必须记住的三个提醒}

\begin{enumerate}[label=\textbf{\arabic*.}]
  \item \textbf{这不是数学证明。} 全文是物理推导，每一步在物理上高度可信，
        但没有一步构成定理。真正的证明路线（Gaitsgory 等人在 2024 年前后完成的纯代数几何路线）
        在本文之后十几年才出现。不要把本文当作几何 Langlands 的证明来引用。
  \item \textbf{只做非分歧情形。} 作者在引言中明确说：分歧（ramified）情形不在本文范围。
        后来 Witten 2008 与 Frenkel 2009 的部分工作就是往分歧方向扩展。
  \item \textbf{第 12 节是诚实度指标。} 前面所有膜都活在"有效二维理论"里；
        \S12 才把它们装回四维规范理论。因此 \S12 最技术，也最容易暴露前面推导的漏洞。
        只读 \S11 不读 \S12，会误以为整个对应是严格的。
\end{enumerate}

\section*{五、常见的十个误解}

\begin{enumerate}[label=\arabic*.,leftmargin=2.2em]
  \item \emph{"KW 证明了几何 Langlands。"} 没有。他们给出的是物理推导与一致性检查。
  \item \emph{"$\Psi$ 就是耦合常数。"} 不是。$\Psi$ 是扭化参数；耦合 $\tau$ 是另一个量
        （\S2 定义），两者在 \S3 的关系是本文的技术核心之一。
  \item \emph{"$S$-对偶交换 $G$ 与 $\Ldual$。"} 严格说交换的是规范群与其 Langlands 对偶群，
        在 \S2 用 GNO 的磁对偶群来钉死。写法上要注意原文用的是 $\Ldual$ 而不是 $\hat G$。
  \item \emph{"膜就是 D-brane 的复数。"} 本文的膜是二维 $\sigma$-模型上的 A-膜或 B-膜，
        定义依赖三个复结构 $I,J,K$ 中的哪一个被当作"复结构"、哪一个被当作"辛结构"。
  \item \emph{"特征膜只是特征层的物理名字。"} 本征性有精确含义：线算子作用在膜上等于乘以一个数（本征值）。
        这个"数"才是 Hecke 特征值。
  \item \emph{"Hecke 算子只在 \S9 出现。"} 它在 \S10 被"算"了一遍；只看 \S9 会以为它是定义，实则不然。
  \item \emph{"c.c.\ A-膜是一个普通的 A-膜。"} 它是余各向同性的（coisotropic），不是拉格朗日的；
        这一点在 \S11.1 专门交代。
  \item \emph{"\S7 可以完全跳过。"} 第一遍可以；但若关心群中心的离散通量与 $\pi_1(G)$ 的相位修正，
        \S7 是唯一出处。
  \item \emph{"镜像对称就是 T-对偶。"} 在本文语境里是 $\MH${} 上 $J\leftrightarrow K$ 的交换，
        与 T-对偶不是同一件事。
  \item \emph{"读完就会做研究了。"} 本文是入口，不是终点。读完应能接着读 Witten 2008 与 Frenkel 2009。
\end{enumerate}

\section*{六、和两篇必读综述的关系}

\textbf{Witten 2008}（\emph{Mirror Symmetry, Hitchin's Equations, And Langlands Duality}，
arXiv:0802.0999，13 页）是本文的\textbf{短版重述}：作者用更几何的语言把同一个故事讲了一遍，
删掉了大部分物理推导，突出了镜像对称与 Hitchin 方程的位置。读完 \kw{} 的 \S5 之后读它最省力。

\textbf{Frenkel 2009}（\emph{Gauge Theory And Langlands Duality}，arXiv:0906.2747，32 页，
Séminaire Bourbaki 报告）是\textbf{数学家的导读}：它从几何 Langlands 的角度复述 \kw{} 的结论，
指出哪些是物理推导、哪些在数学上还需要补，并讨论了后续进展。读者若想判断
"哪些结论今天已经严格化"，应当以它为准。

两篇的关系可以这样记：\kw{} 是\textbf{原始论文}，Witten 2008 是\textbf{作者自己的摘要}，
Frenkel 2009 是\textbf{数学家写的审稿意见}。

\section*{七、译本体例速查}

\begin{itemize}[leftmargin=1.6em]
  \item 正文之前每节有\textbf{本节重点}（蓝框）：3--5 条，给出物理图像 / 数学工具 / 关键结论 / 边界。
  \item 每片末尾有\textbf{译注}（橙框）：易错点、含糊之处、物理图像补充、原文疑误。
        补白一律以「按：」起首，并标来源等级：\textbf{[A]} 原文明确、\textbf{[B]} 权威综述或百科、
        \textbf{[I]} 译者推断。
  \item 公式编号与原文一一对应；译文里保留了 \texttt{\% 式 (2.15)} 形式的锚点注释，方便对照 PDF。
  \item 原文笔误照录，另加「[原文疑误]」译注说明。
  \item 页锚：每节开头给出原文印刷页范围，脚注形如 \src{源 p.56}。
\end{itemize}


\section*{八、逐节导览：每一节该怎么读}

下面按原文 12 节各给一段导读。\textbf{目标}是让你在翻开某一节之前知道"这节要交什么货"，
读完之后知道"我该不该读懂每一个细节"。

\paragraph{\S1 引言（源 pp.2--5）}
四页纸塞进了全篇的六个基本想法：扭化、$\Nfour$ 对偶给出 Langlands 对偶、
紧化到曲线、$\MH${} 上的膜、't Hooft 算子与 Hecke 算子、A-膜与 $\mathcal{D}$-模。
\textbf{第一次读几乎必然会觉得抽象}，这是正常的：引言的写法是"先给你结论，再在正文里兑现"。
建议先只记住三句话：(a) 四维 $S$-对偶的二维化身是镜像对称；(b) 线算子作用在膜上，
特征膜就是 Hecke 特征层；(c) 作者明确声明不含分歧情形。其余回头再看。

\paragraph{\S2 $\Nfour$ 超对称 Yang--Mills 与 $S$-对偶（源 pp.6--15）}
这节是全篇的\textbf{词表}。要拿到的东西只有三样：$\Nfour$ 理论的场内容（规范场、四个 Weyl 费米子、
六个标量）、复耦合 $\tau=\theta/2\pi+4\pi i/g^2$、以及 $S$-对偶 $\tau\mapsto-1/\tau$。
真正的关键是 \textbf{GNO 的磁对偶群恰好就是 Langlands 对偶群}：一旦承认这一点，
"电磁对偶"和"Langlands 对偶"就从两个领域的词变成同一个词。Table~1 给出 $G$ 与 $\Ldual$ 的对应表
（$SU(N)\leftrightarrow SU(N)$、$SO(2n+1)\leftrightarrow Sp(2n)$ 等），值得抄在纸上。

\paragraph{\S3 由 $\Nfour$ SYM 得到拓扑场论（源 pp.16--37）}
这一节最技术，也最容易把人劝退：它要做的是"扭化"（twist），即把超对称荷重新组合成标量荷，
使理论变成拓扑场论。\textbf{读懂的标准不是会算，而是接受三件事}：(a) $\Nfour$ 允许一整族扭化；
(b) 这一族由单个复参数 $\Psi$ 参数化；(c) 特殊值 $\Psi=0,\infty,\pm1,\pm i$ 对应已知的拓扑理论。
\textbf{卡点}：原文在这里用 $\mathcal{Q}$ 记 BRST 荷、用 $\delta$ 记其作用，符号密集；第一遍可以
只读每小节开头与"特殊值表"，细节留到需要用的时候再回来。

\paragraph{\S4 紧化与 $\MH$ 的几何（源 pp.38--55）}
把四维理论紧化在曲线 $C$ 上（$C$ 是紧黎曼面，亏格 $g$），二维有效理论的目标空间就是
\textbf{Hitchin 模空间} $\MH$：$C$ 上 $G$-Higgs 对的模空间。它是一个\textbf{超凯勒流形}，
所以有三个复结构 $I,J,K$（满足四元数关系）。\textbf{必须记住}：$\MH$ 作为实流形的维数是
$4(g-1)\dim G$，并且它带有 Hitchin 纤维化，纤维是阿贝尔簇（Jacobi 或 Prym）。
\S4 还说明这个系统是\textbf{完全可积}的——后面 \S9--\S10 的一切都建立在这个几何上。

\paragraph{\S5 二维拓扑场论（源 pp.56--77）}
\textbf{全篇第一个核心。}在 $\MH${} 上做二维 $\sigma$-模型，取三个复结构中的不同组合就得到 A-模型或 B-模型。
$S$-对偶在四维交换 $G$ 与 $\Ldual$；到了二维，它表现为把 $J$ 与 $K$ 的角色对调，
也就是\textbf{镜像对称} $B(J)\cong A(K)$。这一节还引入广义复几何的语言，用来统一描述
"哪个复结构被当作复结构、哪个被当作辛结构"。\textbf{读法}：先接受 $S$-对偶 $=$ 镜像对称这个结论，
再回看 \S3 的特殊值表，会发现 $\Psi$ 的取值正好在挑选用哪一对复结构。

\paragraph{\S6 圈算子与线算子（源 pp.78--102）}
四维理论里的 Wilson 线与 't Hooft 线，紧化到 $C$ 上之后变成二维理论里的\textbf{算子}（插入在点上）。
这节把"算子作用在态上"翻译成"算子作用在膜上"，并定义\textbf{特征膜}：
若某线算子作用在膜 $\mathcal{B}$ 上等于乘一个数，则称 $\mathcal{B}$ 是特征膜。
\textbf{这是整个对应的语法核心}——后面所有"$\mathcal{D}$-模 / Hecke 特征层"都是特征膜的具体化。

\paragraph{\S7 通量与 $S$-对偶（源 pp.103--115）}
群中心 $Z(G)$ 的离散通量（'t Hooft 通量）会在 $C$ 上留下相位，使 $S$-对偶的陈述需要精细修正。
\textbf{第一遍可以整节跳过}，只要知道"存在这些离散修正，它们细化但不推翻 \S5 的结论"。
若日后要做带通量的对应（例如涉及 $\mathrm{Bun}_G$ 的连通分支），再回来读。

\paragraph{\S8 电特征膜（源 pp.116--122）}
最短的一节，但地位重要：它说明电的 Wilson 线的特征膜与 $\pi_1(C)\to\LdualC$ 的表示
（即 $C$ 上的局部系统、平坦联络）一一对应。数学上这就是几何 Langlands 的
"自守侧 $\leftrightarrow$ 局部系统"在物理里的化身。\textbf{读懂标准}：能说清
"为什么 Wilson 线的本征值给出一个表示"。

\paragraph{\S9 't Hooft 算子与 Hecke 算子（源 pp.123--147）}
\textbf{全篇第二个核心。}把 't Hooft 算子插在 $C$ 的一点上，它对膜的作用等价于
在对应的 $\mathcal{D}$-模上做一次 \textbf{Hecke 修改}（在一点上改变丛的结构，保留其余部分）。
Hitchin 模空间在这里与仿射 Grassmann 流形发生联系：Hecke 修改的模空间就是仿射 Grassmann 的
某个 Schubert 胞腔。\textbf{读法}：\S9 给出"同一性"的陈述，真正算出来在 \S10，所以
读 \S9 时不要急着追问"为什么等于"，先记住对应关系。

\paragraph{\S10 Bogomolny 方程与 Hecke 修改（源 pp.148--183）}
最长的技术节，任务是把 \S9 的同一性\textbf{算}出来：用\textbf{单极子冒泡}（monopole bubbling）
实现 Hecke 修改。技术链条是：边界条件 $\to$ 单极子解 $\to$ Kähler 商 $\to$ 模空间同构。
\textbf{卡点有两个}：\S10.4 的算子乘积展开（OPE）与 \S10.5 的扩展 Bogomolny 方程。
\textbf{建议}：第一遍只读 \S10.1--10.2 的动机与结论，把 OPE 与扩展方程留给第二遍。

\paragraph{\S11 A-膜与 $\mathcal{D}$-模（源 pp.184--206）}
\textbf{收口。}要把 A-膜（辛几何对象）翻译成 $\mathcal{D}$-模（代数几何对象），需要一个"翻译器"：
\textbf{典范余各向同性 A-膜} $\mathcal{B}_{cc}$。它不是拉格朗日膜，而是余各向同性的——
这一节的 \S11.1 专门交代它的定义（余各向同性子流形 = 若干 Poisson 交换函数的公共零点集）。
然后作者证明：$\mathcal{B}_{cc}$ 与任意 A-膜 $\mathcal{B}$ 之间的态空间
$\mathrm{Hom}(\mathcal{B}_{cc},\mathcal{B})$ 给出 $\mathcal{D}$-模，且对 $\mathcal{B}$ 的函子性成立。
\textbf{读法}：这一节是代数几何读者最该细读的一节，因为它是唯一把物理论证翻译成
可检验数学陈述的地方。

\paragraph{\S12 由规范理论得到膜（源 pp.207--222）}
\textbf{诚实度指标。}前面所有膜都活在"有效二维理论"（$\sigma$-模型）里，\S12 才把它们
提升回四维规范理论的边界条件，得到三类膜：$(B,B,B)$、$(B,A,A)$、$(A,B,A)$。
读者若想知道"哪些推导依赖有效理论近似"，答案全在这一节。
\textbf{建议}：第二遍再读；第一遍只需知道它存在，并记住作者的结论——
物理图像自洽，但四维层面的严格性仍有缺口。

\section*{九、符号速查}

\begin{center}\small
\begin{tabular}{@{}ll@{\hspace{2em}}ll@{}}
\toprule
符号 & 含义 & 符号 & 含义 \\
\midrule
$C$ & 紧黎曼面（曲线） & $G$ & 紧单连通规范群 \\
$\Ldual$ & Langlands 对偶群 & $\LdualC$ & 其对偶复群 \\
$\tau$ & 复耦合常数 & $\Psi$ & 扭化参数（典范参数）\\
$\MH$ & Hitchin 模空间 & $I,J,K$ & 三个复结构 \\
$\mathcal{M}_{\mathrm{flat}}$ & 平坦联络模空间 & $B(J),A(K)$ & B-膜、A-膜 \\
$\mathcal{B}_{cc}$ & 典范余各向同性 A-膜 & $\mathcal{B}_{op}$ & oper 之膜 \\
$W_R(\gamma)$ & Wilson 线 & $T(\gamma)$ & 't Hooft 线 \\
$\mathrm{Bun}_G$ & $G$-丛模空间 & $\mathrm{Hecke}$ & Hecke 对应 \\
\bottomrule
\end{tabular}
\end{center}

\vfill
\begin{center}\small\color{gray}
（本导读为译本附册，供首次阅读原文者使用；不替代原文，也不构成任何数学证明。）
\end{center}

\end{document}
